%CLASSE DOCUMENTO - LINGUA E DIMENSIONE FONT
\documentclass[11pt,english]{toptesi}

%%%%%%%%%%%%%%%%%%%%%%%%%%%%%%%%%%%%%%%%%%%%%%%%%%%%%%%%%%%%%%%

% INCLUSIONE PACCHETTI
\usepackage[utf8]{inputenc} %utf8
\usepackage[english]{babel}
\usepackage[T1]{fontenc}
\usepackage{blindtext}
\usepackage{graphicx,wrapfig}
\usepackage{booktabs}
\usepackage{lmodern}
\usepackage{varioref}
\usepackage{url}
\usepackage{array}
\usepackage{paralist}{\obeyspaces\global\let =\space}
\usepackage{verbatim}
\usepackage{subfig}
\usepackage{tabularx}
\usepackage{amsmath}
\usepackage{amsfonts}
\usepackage{float}
\usepackage{amssymb}
\usepackage{multicol}
\usepackage{multirow}
\usepackage{listings}
\usepackage[pass]{geometry}
\usepackage[figuresright]{rotating}
\usepackage{algorithm}
\usepackage{algorithmic}
\usepackage{amsmath}
\usepackage[babel]{csquotes}
\usepackage{hyperref}
\usepackage[backend=bibtex]{biblatex}
\usepackage{tikz}

%%%%%%%%%%%%%%%%%%%%%%%%%%%%%%%%%%%%%%%%%%%%%%%%%%%%%%%%%%%%%%%

% CONFIGURAZIONE LINK E RIFERIMENTI
\hypersetup{%
    pdfpagemode={UseOutlines},
    bookmarksopen,
    pdfstartview={FitH},
    colorlinks,
    linkcolor={black}, %COLORE DEI RIFERIMENTI AL TESTO
    citecolor={blue}, %COLORE DEI RIFERIMENTI ALLE CITAZIONI
    urlcolor={blue} %COLORI DEGLI URL
}

%%%%%%%%%%%%%%%%%%%%%%%%%%%%%%%%%%%%%%%%%%%%%%%%%%%%%%%%%%%%%%%

% CONFIGURAZIONE LISTATI/CODICE - CANCELLARE SE NON NECESSARIO
% PYTHON - BIANCO E NERO
\lstset{%
	captionpos=b,
	language=Python,
	basicstyle =\small\ttfamily,
	keywordstyle=\color{black}\bfseries,
	breaklines=true,
	breakatwhitespace=true,
	frame=lines,
	numbers=left,
	numberstyle=\footnotesize,
}

%%%%%%%%%%%%%%%%%%%%%%%%%%%%%%%%%%%%%%%%%%%%%%%%%%%%%%%%%%%%%%%

% FRENCHSPACING ABILITATO - CANCELLARE PER SPAZIATURA ALL'INGLESE
%\frenchspacing

%%%%%%%%%%%%%%%%%%%%%%%%%%%%%%%%%%%%%%%%%%%%%%%%%%%%%%%%%%%%%%%

%DEFINIZIONE SEZIONI IN NUMERAZIONE ROMANA
%ELENCO DEI LISTATI/CODICI
\makeatletter
\newcommand\listofcodes{%
 \iffrontmatter\else\frontmattertrue\fi
 \if@openright\cleardoublepage\else\clearpage\fi
 % change the meaning of \chapter in a group
 \begingroup\def\chapter##1{\@schapter}
 \phantomsection % for the hyperlink
 \lstlistoflistings
 \endgroup
}
\makeatother

%%%%%%%%%%%%%%%%%%%%%%%%%%%%%%%%%%%%%%%%%%%%%%%%%%%%%%%%%%%%%%%

% INFORMAZIONI PDF - PERSONALIZZARE
\pdfinfo{%
  /Title    (Control oriented modelling of four wheeled electric vehicles)
  /Author   (Sebastian Giles)
  /Subject  (How to write meaningful titles)
  /Keywords (LaTeXi foo bar baz meaningful PIT Presicce)
}

%%%%%%%%%%%%%%%%%%%%%%%%%%%%%%%%%%%%%%%%%%%%%%%%%%%%%%%%%%%%%%%

% FRONTESPIZIO - PERSONALIZZARE
% ELIMINATE LE VOCI CHE NON VI SERVONO

% LOGO UNIVERSITA
\logosede{images/logo}

% UNIVERSITA - NOME
\ateneo{Università Politecnica delle Marche}

% FACOLTA - DICITURA
\FacoltaDi{Faculty of }
% FACOLTA - NOME
\facolta{Engineering}

% CORSO DI LAUREA - DICITURA (MANTENERE LO SPAZIO)
\CorsoDiLaureaIn{Bachelor of Science in }
% CORSO DI LAUREA - NOME
\corsodilaurea{Computer and Automation Engineering}

% TIPOLOGIA TESI
\TesiDiLaurea{Bachelor's Thesis}

% TITOLO
\titolo{Control oriented modelling of four wheeled electric vehicles}

% SOTTOTITOLO
%\sottotitolo{Bella raga}

% RELATORE/I - DICITURA
\AdvisorName{Advisor}
% RELATORE - PROF. NOME E COGNOME
\relatore{prof.\ Andrea Bonci}
% RELATORE AGGIUNTIVO - PROF NOME E COGNOME
% SE SI HA SOLO UN RELATORE ELIMINARE E CAMBIARE Advisors in Advisor
%\secondorelatore{prof.\ Neo Cortex}

% TUTORE AZIENDALE - TITOLO NOME E COGNOME
%\tutoreaziendale{Ing. Pug Dog}
% TUTORE AZIENDALE - DICITURA//AZIENDA
%\NomeTutoreAziendale{Company tutors\\FeelGood Inc}

% CANDIDATO - DICITURA (MANTENERE I DUE PUNTI)
\CandidateName{Candidate:}
% SECONDO CANDIDATO - ELIMINARE O DECOMMENTARE
%secondocandidato{Bombo de Bombis}

% CANDIDATO - NOME E COGNOME
\candidato{Sebastian Giles}


% DATA - MESE ANNO
\sedutadilaurea{July 2018}

%%%%%%%%%%%%%%%%%%%%%%%%%%%%%%%%%%%%%%%%%%%%%%%%%%%%%%%%%%%%%%%

% LISTA DEI CAPITOLI DA INCLUDERE - PERSONALIZZARE
% \includeonly{%
% 01-intro,%
% 02-stuff,%
% 03-conclusion,%
% 04-appendix-A,%
% }

% FILE DI BIBLIOGRAFIA
\bibliography{bibliography}


%%%%%%%%%%%%%%%%%%%%%%%%%%%%%%%%%%%%%%%%%%%%%%%%%%%%%%%%%%%%%%%

% INIZIO DOCUMENTO
\begin{document}
\english

\frontespizio

%%%%%%%%%%%%%%%%%%%%%%%%%%%%%%%%%%%%%%%%%%%%%%%%%%%%%%%%%%%%%%%

%INTERLINEA - DEFAULT 1 - NON ESAGERATE, NON SUPERATE MAI 1.3 ;)
\interlinea{1.3}

%%%%%%%%%%%%%%%%%%%%%%%%%%%%%%%%%%%%%%%%%%%%%%%%%%%%%%%%%%%%%%%

\frontmatter

% DEDICA - PERSONALIZZARE
% VSPACE - PROPORZIONE USATA PER CENTRATURA VERTICALE DEL TESTO
% FLUSHRIGHT - ALLINEAMENTO ORIZZONTALE A DESTRA
\vspace*{\stretch{1}}
\begin{flushright}
\noindent
To me and only me
\end{flushright}
\vspace*{\stretch{6}}
\cleardoublepage

% CITAZIONE - PERSONALIZZARE
% VSPACE - PROPORZIONE USATA PER CENTRATURA VERTICALE DEL TESTO
% FLUSHRIGHT - ALLINEAMENTO ORIZZONTALE A DESTRA
\vspace*{\stretch{1}}
\begin{flushright}
\noindent
Waka waka. Eh-eh.

\textit{Shakira}
\end{flushright}
\vspace*{\stretch{6}}
\cleardoublepage

%%%%%%%%%%%%%%%%%%%%%%%%%%%%%%%%%%%%%%%%%%%%%%%%%%%%%%%%%%%%%%%

% RINGRAZIAMENTI - PERSONALIZZARE
\ringraziamenti
Thanks to someone

%%%%%%%%%%%%%%%%%%%%%%%%%%%%%%%%%%%%%%%%%%%%%%%%%%%%%%%%%%%%%%%

% ABSTRACT - PERSONALIZZARE
\sommario
Abstract, pleased to meet you. Ah, it's not funny anymore.

%%%%%%%%%%%%%%%%%%%%%%%%%%%%%%%%%%%%%%%%%%%%%%%%%%%%%%%%%%%%%%%

% INDICI - ELIMINARE GLI INDICI NON NECESSARI

% INDICE GENERALE
\tableofcontents

% INDICE DELLE FIGURE
\listoffigures

% INDICE DELLE TABELLE
\listoftables

% INDICE DEI CODICI
\listofcodes

%%%%%%%%%%%%%%%%%%%%%%%%%%%%%%%%%%%%%%%%%%%%%%%%%%%%%%%%%%%%%%%

\mainmatter

% INCLUSIONE FILE CAPITOLI - PERSONALIZZARE - TENERE COERENTE CON LISTA IN ALTO
\chapter{Introduction}
\label{chap:intro}
\section{Background}
\label{sec:background}

\section{Objectives}
\label{sec:objectives}

\section{Approach}
\label{sec:approach}

\chapter{Qui}
\label{chap:qui}

\section{Quaraquaraqua}
\label{sec:quaqaraqua}

This is a reference to a chapter \ref{chap:quo}. This is a reference to a figure \ref{fig:doge}. This is a reference to some code \ref{lst:hello}. This is a citation \cite{famous:paper}.

\lstinputlisting[label=lst:hello, firstline=2, lastline=4, caption={I directly included a portion of a file}]{code/hello.py}

\begin{lstlisting}[language=Java, label=lst:java, caption={Some code in another language than the default one}]
public void prepare(AClass foo) {
  AnotherClass bar = new AnotherClass(foo)
}
\end{lstlisting}

\Blindtext


%    \usepackage{fullpage}
%    \usepackage{fourier}

\usetikzlibrary{calc,3d}

\thispagestyle{empty}

\begin{center}
  \begin{tikzpicture}[x  = {(-0.5cm,-0.5cm)},
    y  = {(0.9659cm,-0.25882cm)},
    z  = {(0cm,1cm)},
    scale = 2,
    color = {lightgray}]
    % style of faces
    \tikzset{facestyle/.style={fill=lightgray,draw=red,very thin,line join=round}}
    % face "back"
    \begin{scope}[canvas is zy plane at x=0]
      \path[facestyle,shade] (0,0) rectangle (2,4);
    \end{scope}
    % face  "left"
    \begin{scope}[canvas is zx plane at y=0]
      \path[facestyle,shade] (0,0) rectangle (2,3);
    \end{scope}
    % face "front"
    \begin{scope}[canvas is zy plane at x=3]
      \path[facestyle] (0,0) rectangle (2,4);
    \end{scope}
    % face  "right"
    \begin{scope}[canvas is zx plane at y=4]
      \path[facestyle] (0,0) rectangle (2,3);
    \end{scope}
    % face "up"
    \draw[fill=brown,draw=red,opacity=.8,very thin,line join=round]
    (0,0,2) --
    (3,0,2) --
    (3,{4*cos(15)},{4*sin(15)+2}) --
    (0,{4*cos(15)},{4*sin(15)+2}) --cycle ;
    % labels
    \draw[very thin,black,line join=round]
    (3,0,0) -- node [sloped,below] {$a$}
    (3,4,0) -- node [right]        {$b $}
    (3,4,2) -- node [sloped,above] {$c=\sqrt{a^2+b^2}$}
    (3,0,0);
  \end{tikzpicture}
\end{center}

\begin{figure}
  \begin{center}
    \includegraphics[width=0.5\columnwidth]{images/doge.png}
  \end{center}
  \caption{This is not a figure. It's a caption.}
  \label{fig:doge}
\end{figure}

\chapter{Quo}
\label{chap:quo}
\Blindtext

\appendix
% INCLUSIONE APPENDICI - - PERSONALIZZARE - TENERE COERENTE CON LISTA IN ALTO
\chapter{An appendix}
\label{app:a}
\Blindtext


%%%%%%%%%%%%%%%%%%%%%%%%%%%%%%%%%%%%%%%%%%%%%%%%%%%%%%%%%%%%%%%

% BIBLIOGRAFIA
\addcontentsline{toc}{chapter}{\refname}
\nocite{*}
\printbibliography

\end{document}
